\section{Canonical Multi-Secant Recovery}
\label{sec:multisecant-recovery}

The posterior-certified exact solver computes one canonical action projection.  Nested action
fibers turn those projections into a finite inverse-shape recovery law.  Let
\(H\in\SPD^d\) be a fixed quadratic Hessian and define its determinant-one
inverse shape
\begin{equation}
  P_\star=c_HH^{-1},
  \qquad c_H=(\det H)^{1/d}.
  \label{eq:inverse-shape}
\end{equation}
The scalar \(c_H\) is assumed known, or supplied by an independent scalar
gauge.  This information requirement is part of the theorem.

Let \(S^{(1)},\ldots,S^{(K)}\) have strictly nested column spaces, delete
dependent columns, and write \(r_k=\rank S^{(k)}\).  Define the cumulative
shape-normalized actions
\begin{equation}
  A_k=HS^{(k)},
  \qquad B_k=c_HS^{(k)},
  \qquad
  \mathscr C_k=\{P\in\Pdet:PA_k=B_k\}.
  \label{eq:nested-action-fibers}
\end{equation}
Then \(P_\star\in\mathscr C_k\) and
\(\mathscr C_{k+1}\subset\mathscr C_k\).  Starting from any
\(P_0\in\Pdet\), define
\begin{equation}
  P_k=\argminop_{P\in\mathscr C_k}\dai(P_{k-1},P).
  \label{eq:canonical-multisecant-update}
\end{equation}

\begin{theorem}[Nested canonical multi-secant identification]
\label{thm:nested-multisecant}
For the exact outer projections in \eqref{eq:canonical-multisecant-update},
every update exists uniquely and
preserves all earlier cumulative secants.  For every \(k\),
\begin{equation}
  \boxed{
  \dai(P_k,P_\star)^2+
  \dai(P_{k-1},P_k)^2
  \le\dai(P_{k-1},P_\star)^2.}
  \label{eq:multisecant-pythagorean}
\end{equation}
Consequently,
\begin{align}
  \sum_{j=1}^k\dai(P_{j-1},P_j)^2
  &\le\dai(P_0,P_\star)^2-
       \dai(P_k,P_\star)^2,
  \label{eq:multisecant-energy}\\
  \sum_{j=1}^k\dai(P_{j-1},P_j)
  &\le\sqrt{k}\,\dai(P_0,P_\star).
  \label{eq:multisecant-length}
\end{align}
The determinant-one identification threshold is sharp:
\begin{equation}
  r_K=d-1\quad\Longrightarrow\quad P_K=P_\star,
  \label{eq:sharp-identification-sufficiency}
\end{equation}
whereas any cumulative rank \(r\le d-2\) leaves infinitely many feasible
determinant-one SPD completions.  If each pre-final round contributes exactly
\(b\) new independent probe directions and the final round contributes the
remainder, exact recovery occurs after
\begin{equation}
  K=\left\lceil\frac{d-1}{b}\right\rceil
  \label{eq:block-identification-rounds}
\end{equation}
rounds.  Each exact outer projection reduces by AIRM congruence to the
identity-reference problem in \cref{thm:certified-action-solver}; the inner
iteration converges globally to it and supplies a residual certificate for a
finite approximation.
\end{theorem}

The outer decrease remains quantitative when each projection is stopped after
finitely many certified inner iterations.  Set \(\widetilde P_0=P_0\).  At
round \(k\), let
\[
  Q_k=\operatorname{Proj}_{\mathscr C_k}(\widetilde P_{k-1})
\]
be the exact projection used only for analysis, and let
\(\widetilde P_k\in\mathscr C_k\) be the feasible inner output.  This theorem
uses the exact-feasibility model: action-chart inner iterates lie in
\(\mathscr C_k\) in exact arithmetic.  Floating-point feasibility error is a
separate perturbation and is not included in \(\varepsilon_k\).

\begin{theorem}[Residual-controlled finite inner projections]
\label{thm:approximate-multisecant}
Suppose a posterior inner certificate gives
\begin{equation}
  \dai(\widetilde P_k,Q_k)\le\varepsilon_k.
  \label{eq:inner-projection-certificate}
\end{equation}
Then, with \(D_k=\dai(\widetilde P_k,P_\star)\),
\begin{align}
  &\dai(\widetilde P_{k-1},\widetilde P_k)^2+D_k^2
  \nonumber\\
  &\qquad\le D_{k-1}^2
  +2\sqrt2\,\varepsilon_kD_{k-1}+2\varepsilon_k^2,
  \label{eq:approximate-multisecant-pythagorean}\\
  D_k&\le D_{k-1}+\varepsilon_k
  \le D_0+\sum_{j=1}^k\varepsilon_j.
  \label{eq:approximate-multisecant-distance}
\end{align}
Consequently,
\begin{equation}
  \sum_{j=1}^k
  \dai(\widetilde P_{j-1},\widetilde P_j)^2
  \le D_0^2-D_k^2
  +\sum_{j=1}^k
   \left(2\sqrt2\,\varepsilon_jD_{j-1}+2\varepsilon_j^2\right).
  \label{eq:approximate-multisecant-energy}
\end{equation}
All earlier cumulative actions remain exact because
\(\widetilde P_k\in\mathscr C_k\subseteq\mathscr C_j\) for \(j<k\).
After congruence to the identity-reference inner problem,
\cref{eq:solver-distance-certificate} permits the observable choice
\(\varepsilon_k=\fro{W_{B,k}(\widetilde\Sigma_k)}\); under the residual-oracle
model, \cref{eq:inexact-certificates} supplies the corresponding enlarged
bound.  If \(r_K=d-1\), then \(\mathscr C_K=\{P_\star\}\), so every feasible
inner output at the final round equals \(P_\star\) exactly.
\end{theorem}

The known scalar gauge cannot be dropped.  Ordinary Hessian-vector secants in
\(d-1\) directions do not determine the inverse shape.  In dimension two,
\[
  H_t=\operatorname{diag}(1,t),\qquad s=e_1
\]
gives the same observation \(H_ts=e_1\) for every \(t>0\), while
\[
  (\det H_t)^{1/2}H_t^{-1}
  =\operatorname{diag}(\sqrt t,1/\sqrt t)
\]
varies with \(t\).  Thus \eqref{eq:sharp-identification-sufficiency} is a
sharp theorem for shape-normalized exact actions, rather than a claim of
scale-free identification from ordinary secants.  Acquisition of \(c_H\) is
external to the action oracle and is not claimed to be cheaper than an
additional probe.
